\documentclass[10pt,a4paper]{amsart}
\usepackage[margin=19mm]{geometry}
\usepackage{amsmath,amssymb,mathtools}
\usepackage[scr=rsfs,cal=euler]{mathalfa}
\usepackage{xfrac}
\usepackage[unicode,hidelinks,bookmarks=false]{hyperref}
\newcommand{\ie}{i.e.\ }
\newcommand{\cf}{cf.\ }
\newcommand{\resp}{respectively\ }
\newcommand{\Subsection}{\subsection}
\providecommand{\nobreakdash}{\nobreak\hbox{-}}
\setlength{\emergencystretch}{2em}
\multlinegap0pt
\usepackage{mathtools}
\usepackage[shortlabels]{enumitem}
\usepackage{tikz}
\usepackage{amssymb}
\usepackage{xcolor}
\usepackage[normalem]{ulem}
\usepackage{algorithm}\usepackage{algpseudocode}
\newtheorem{proposition}{Proposition}
\newcommand{\Ber}{\mathsf{Ber}}
\title{Breaking the Finite-Sample Barrier in Entropy Coupling}
\author{Shahab Asoodeh, Jun Chen}
\date{Source: arXiv:2605.16229v1 (CC BY 4.0)}
\begin{document}
\maketitle

\noindent\emph{Source and licence.} Breaking the Finite-Sample Barrier in Entropy Coupling. Version arXiv:2605.16229v1 (CC BY 4.0), distributed by arXiv under the Creative Commons Attribution 4.0 International licence, http://creativecommons.org/licenses/by/4.0/, as recorded in the arXiv licence metadata for this exact version and shown on the abstract page https://arxiv.org/abs/2605.16229v1. Full licence notice: \texttt{licenses/arxiv-2605.16229-CC-BY-4.0.txt}.\par\smallskip

\noindent\textit{Editorial scope.} Counted unit: the article's proposition Prop:Binary\_m (source0.tex lines 447--464). The article's complete original proof of it is delivered with it (source0.tex lines 1155--1382, one \texttt{proof} environment of 228 lines, which the article places away from the statement). No mathematics is written, repaired or re-derived: the statement and the proof are the article's own lines, in the article's own notation, and no retained line is substituted or re-flowed.  No further source-local context is needed: every cross-reference of every retained line resolves inside the two delivered spans.  Nothing was omitted from any retained span, so the package carries no omission record.  No retained line contains a citation, so the wrapper carries no bibliography and no bibliography entry.  No retained line contains a figure environment, an image-inclusion command or drawing code, so no image is delivered and the package declares no asset dependency.  Editorial formatting only, in addition: an AMS \texttt{amsart} preamble that restates the article's own declarations and macros used by the retained lines, namely the environments \texttt{proposition} on separate counters, together with the standard packages those lines need.  The retained environments print in this document as Proposition 1 in order of appearance, whereas the article prints the same items with the numbering of its own full document.  The complete native source of the article as distributed in the arXiv e-print for this exact version is bundled unchanged as \texttt{sources/200712/source0.tex}.
\begin{proposition}\label{Prop:Binary_m}
For $m\ge2$, we have
$$
H_m(\Ber(p)\|\Ber(q)) =
\begin{cases}
0, & r \le m s,\\
(1-ms)\,h_{\mathsf b}\!\big(\frac{r-ms}{1-ms}\big), & r>ms,
\end{cases}
$$
where $r=\min\{p,1-p\}$ and $s=\min\{q,1-q\}$. For $m=1$, we have
$$
H_1(\Ber(p)\|\Ber(q)) =
\begin{cases}
s\,h_{\mathsf b}\!\big(\frac{r}{s}\big), & r\le s,\\
(1-s)\,h_{\mathsf b}\!\big(\frac{r-s}{1-s}\big), & r>s.
\end{cases}
$$
\end{proposition}

\begin{proof}[Proof of Proposition~\ref{Prop:Binary_m}]
We first give the proof for $m\geq 2$ and then for $m=1$. 
Notice that  
\[H_m(\Ber(p)\|\Ber(q)) =H_m(\Ber(1-p)\|\Ber(q)),\]
which is obtained by replacing $X$ with $1-X$, and
\[H_m(\Ber(p)\|\Ber(q))=H_m(\Ber(p)\|\Ber(1-q)),\]
obtained by replacing each $Y_i$ with $1-Y_i$. Hence it suffices to prove theorem for \(0\le p\le \tfrac12,\) and \(0\le q\le \tfrac12.\)



\noindent\underline{$\bf m\geq 2$:}
Under these assumptions, $r=p$ and $s=q$, so we must show
\[H_m(\Ber(p)\|\Ber(q))=\begin{cases}
0, & p\le mq,\\[1ex]
(1-mq)\,h_{\mathsf{b}}\!\left(\dfrac{p-mq}{1-mq}\right), & p>mq,
\end{cases}\]
To prove this result, we need to prove achievability and converse. We break down the achievability in three cases: (1) $0\le p\le q$, (2) $q\leq  p\le mq$, and (3) $p> mq$.

\paragraph{Case 1: Achievability when $0\le p\le q$.}
We construct $Y^m=(Y_1,\dots,Y_m)$ with the correct marginals and make $X$ a deterministic function of $Y^m$.

Let $\mathbf 1=(1,\dots,1)$ and let $\mathcal W_{m-1}$ denote the set of all binary vectors of length $m$ having exactly $m-1$ ones. Define a probability distribution on $\{0,1\}^m$ by
\[\Pr(Y^m=\mathbf 1)=p,\qquad \Pr(Y^m=\mathbf 0)=1-p-\frac{m(q-p)}{m-1},\]
and for $y\in\mathcal W_{m-1}$
\[\Pr(Y^m=y)=\frac{m(q-p)}{m-1}\cdot \frac1m,\]
and zero elsewhere.
Since $0\le p\le q\le \frac12$ and $m\ge 2$, all masses are nonnegative, indeed
\[1-p-\frac{m(q-p)}{m-1}=1-\frac{mq-p}{m-1} \ge 1-\frac{m/2}{m-1}\ge 0.\]
For each coordinate $j$, we have \(\Pr(Y_j=1)=p+\frac{m-1}{m}\cdot \frac{m(q-p)}{m-1}=q.\) Thus each $Y_j\sim \Ber(q)$.
Now define
\[X:=\mathbf 1\{Y^m=\mathbf 1\}.\]
Then $X$ is a deterministic function of $Y^m$, so \(H(X| Y^m)=0,\)
and \(\Pr(X=1)=\Pr(Y^m=\mathbf 1)=p.\)
Hence $X\sim\Ber(p)$ and
\[H_m(\Ber(p)\|\Ber(q))=0\qquad\text{for }0\le p\le q.\]

\paragraph{Case 2: Achievability when $q\le p\le mq$.}
Let
\[\beta:=\frac{mq-p}{m-1}\in[0,q].\]
Define a distribution on $\{0,1\}^m$ by
\[\Pr(Y^m=\mathbf 1)=\beta,\qquad \Pr(Y^m=\mathbf 0)=1-mq+(m-1)\beta=1-p,\]
\[\Pr(Y^m=e_i)=q-\beta,\qquad i=1,\dots,m,\]
and zero elsewhere, where $e_i$ is the vector with a single $1$ in position $i$.
All masses are nonnegative, and for each coordinate $j$, we have \(\Pr(Y_j=1)=\beta+(q-\beta)=q.\) Thus each $Y_j\sim\Ber(q)$.

Now define
\[X:=\mathbf 1\{Y^m\neq \mathbf 0\}.\]
Then $X$ is a deterministic function of $Y^m$, so \(H(X| Y^m)=0,\)
and \(\Pr(X=1)=\beta+m(q-\beta)=mq-(m-1)\beta=p.\) Hence $X\sim\Ber(p)$ and
\[H_m(\Ber(p)\|\Ber(q))=0
\qquad\text{for }q\le p\le mq.\]

\paragraph{Case 3: achievability when $p>mq$.}
Assume now $p>mq$. Since $p\le \frac12$, this implies $mq<\frac12$.

Define a distribution on $\{0, 1\}^m$ by  \[\Pr(Y^m=\mathbf 0)=1-mq,\qquad \Pr(Y^m=e_i)=q,\qquad i=1,\dots,m,\]
and no mass elsewhere. Then each coordinate marginal is Bernoulli$(q)$.
Given $Y^m$, define $X$ as follows: \(X=1\) deterministically on each \(e_i\) and on the atom $\mathbf 0$ set \(\Pr(X=1| Y^m=\mathbf 0)=\theta,\)
where $\theta\coloneqq \frac{p-mq}{1-mq}.$ Then
\[\Pr(X=1)=mq+(1-mq)\theta=p,\]
so $X\sim\Ber(p)$. Since the only nondeterministic atom is $\mathbf 0$ and 
\[H(X| Y^m)=(1-mq)\,h_{\mathsf{b}}(\theta)=(1-mq)\,h_{\mathsf{b}}\!\left(\frac{p-mq}{1-mq}\right).\]
Therefore
\[H_m(\Ber(p)\|\Ber(q))\le(1-mq)\,h_{\mathsf{b}}\!\left(\frac{p-mq}{1-mq}\right).\]

\paragraph{Converse when $p>mq$.}
Let $(X,Y^m)$ be any coupling with \(X\sim\Ber(p),\) and  \(Y_1,\dots,Y_m\sim\Ber(q).\)
Let
\[t_0:=\Pr(Y^m=\mathbf 0),\qquad x_0:=\Pr(X=1,Y^m=\mathbf 0).\]
Also let \(N:=Y_1+\cdots+Y_m.\)
Then $\mathbb E[N]=mq.$ Since $N\ge \mathbf 1_{\{N>0\}}$, we have  \(\Pr(N>0)\le \mathbb E[N]=mq,\) and hence
\[t_0=\Pr(N=0)\ge 1-mq.\]
Moreover,
\[p=x_0+\Pr(X=1,N>0)\le x_0+\Pr(N>0)\le x_0+mq,\]
implying \(x_0\ge p-mq.\) Now, we can write 
\[H(X| Y^m)=\sum_{y\in\{0,1\}^m}\Pr(Y^m=y)\,
h_{\mathsf{b}}\!\bigl(\Pr(X=1| Y^m=y)\bigr)\ge t_0\,h_{\mathsf{b}}\!\left(\frac{x_0}{t_0}\right).\]
Set \(T:=1-mq\) and \(\delta:=p-mq.\)
Then $t_0\ge T$ and $x_0\ge \delta$. Also, since $p\le \frac12$ and $mq<\frac12$,
\[T> \frac12,\qquad x_0\le p\le \frac12<T.\]
For fixed $x\in[0,\frac12]$, the map \(t\mapsto t\,h_{\mathsf{b}}(x/t)\)
is increasing on $[x,\infty)$, since
\[\frac{\partial}{\partial t}\Bigl[t\,h_{\mathsf{b}}(x/t)\Bigr]=\log\!\frac{t}{t-x}\ge 0,\]
implying 
\[t_0\,h_{\mathsf{b}}\!\left(\frac{x_0}{t_0}\right)\ge
T\,h_{\mathsf{b}}\!\left(\frac{x_0}{T}\right).\]
Now consider
\[f(x):=T\,h_{\mathsf{b}}(x/T),\]
for $x\in[0,T]$. This function is symmetric about $T/2$ and strictly increasing on $[0,T/2]$. Since \(\delta=p-mq\) and \(T-\delta=1-p\ge p,\) every $x\in[\delta,p]$ satisfies \[\min\{x,T-x\}\ge \delta.\]
Also, note that by symmetry and monotonicity, we have
\(f(x)\ge f(\delta),\) for $x\in[\delta,p]$. In particular,
\[T\,h_{\mathsf{b}}\!\left(\frac{x_0}{T}\right)\ge T\,h_{\mathsf{b}}\!\left(\frac{\delta}{T}\right)=(1-mq)\,h_{\mathsf{b}}\!\left(\frac{p-mq}{1-mq}\right).\]
Putting altogether, we can write
\[H(X| Y^m)\ge(1-mq)\,h_{\mathsf{b}}\!\left(\frac{p-mq}{1-mq}\right),\]
which is the desired converse result. 

Now, we prove the theorem for $m=1$.

\noindent\underline{$\bf m= 1$:} 
As before, it suffices to prove the result for \(0\le p\le \tfrac12,\) and \(0\le q\le \tfrac12.\) 

% \textbf{Step 1: Symmetries.}
% We first show that
% \[
% H_1(\Ber(p)\|\Ber(q))
% =
% H(\Ber(1-p)\|\Ber(q))
% =
% H(\Ber(p)\|\Ber(1-q)).
% \]
% Indeed, if $(X,Y)$ is any coupling with $X\sim \Ber(p)$ and $Y\sim \Ber(q)$, then
% \[
% (1-X,Y)
% \]
% is a coupling of $\Ber(1-p)$ and $\Ber(q)$, and since the map $x\mapsto 1-x$ is bijective,
% \[
% H(1-X| Y)=H(X| Y).
% \]
% Taking the infimum over all couplings gives
% \[
% H(\Ber(1-p)\|\Ber(q))
% \le
% H_1(\Ber(p)\|\Ber(q)).
% \]
% Applying the same argument with $p$ replaced by $1-p$ yields the reverse inequality, hence equality. The proof of
% \[
% H_1(\Ber(p)\|\Ber(q))
% =
% H(\Ber(p)\|\Ber(1-q))
% \]
% is identical, replacing $Y$ by $1-Y$.

% Therefore
% \[
% H_1(\Ber(p)\|\Ber(q))
% =
% H(\Ber(r)\|\Ber(s)),
% \]
% where
% \[
% r=\min\{p,1-p\},
% \qquad
% s=\min\{q,1-q\}.
% \]
% Since $r,s\in[0,\frac12]$, it suffices to compute
% \[
% H_1(\Ber(p)\|\Ber(q))
% \qquad\text{under the assumption}\qquad
% 0\le p\le \frac12,\ \ 0\le q\le \frac12.
% \]

Fix $p,q\in[0,\frac12]$. Let $(X,Y)$ be any coupling with
\(X\sim \Ber(p)\) and \(Y\sim \Ber(q).\) Define
\[a:=\Pr(X=1| Y=0),
\qquad b:=\Pr(X=1| Y=1).\]
Then $a,b\in[0,1]$, and since $\Pr(Y=1)=q$ and $\Pr(Y=0)=1-q$, we must have \(\Pr(Y=1)=(1-q)a+qb.\) Similarly, since $X\sim \Ber(p)$, we must have \((1-q)a+qb=p.\)
Conversely, for any $a,b\in[0,1]$ satisfying
\((1-q)a+qb=p,\) if we let $Y\sim \Ber(q)$ and define the conditional law of $X$ by
\[\Pr(X=1| Y=0)=a,\qquad \Pr(X=1| Y=1)=b,\] then $X\sim \Ber(p)$.
Thus the optimization problem is exactly
\[H_1(\Ber(p)\|\Ber(q))
= \min\Bigl\{(1-q)h_{\mathsf{b}}(a)+q\,h_{\mathsf{b}}(b):\ 0\le a,b\le 1,\ (1-q)a+qb=p\Bigr\}.\]
Set
\[F(a,b):=(1-q)h_{\mathsf{b}}(a)+q\,h_{\mathsf{b}}(b).\]
Since $h_{\mathsf{b}}$ is concave on $[0,1]$, $F$ is concave on $[0,1]^2$. The feasible set
\[\mathcal S:=\{(a,b)\in[0,1]^2:(1-q)a+qb=p\}\]
is a compact line segment. A concave function on a compact line segment attains its minimum at an endpoint of the segment. Hence it remains to determine the endpoints of $\mathcal S$ and compare the corresponding values of $F$.

Now we consider two cases. 
\paragraph{Case 1: $p\le q$.}
The line \((1-q)a+qb=p\) meets the sides $a=0$ and $b=0$ of the unit square at
\[(0,p/q) \qquad\text{and}\qquad\Big(\frac{p}{1-q},0\Big),\]
respectively. Since $p\le q$ and $p\le 1-q$, both points lie in $[0,1]^2$. These are the two endpoints of $\mathcal S$. Therefore
\[H_1(\Ber(p)\|\Ber(q))
=\min\left\{
q\,h_{\mathsf{b}}\!\left(\frac{p}{q}\right),
\,(1-q)\,h_{\mathsf{b}}\!\left(\frac{p}{1-q}\right)
\right\}.\]
We now show that the first term is the smaller one. Fix $p\in[0,\frac12]$ and define \(\phi(t):=t\,h_{\mathsf{b}}(p/t)\) for $t\ge p$.
A direct differentiation gives
\[\phi'(t)=\log_2\!\frac{t}{t-p}\ge 0,\]
so $\phi$ is increasing on $[p,\infty)$. Since $q\le 1-q$, we have 
\[q\,h_{\mathsf{b}}\!\left(\frac{p}{q}\right) =\phi(q)\le\phi(1-q)=(1-q)\,h_{\mathsf{b}}\!\left(\frac{p}{1-q}\right).\]
Hence
\[H_1(\Ber(p)\|\Ber(q))=q\,h_{\mathsf{b}}\!\left(\frac{p}{q}\right)
\qquad\text{when } p\le q.\]

\paragraph{Case 2: $p>q$.}
Again, since $p\le \frac12\le 1-q$, the line \((1-q)a+qb=p\)
meets the sides $b=0$ and $b=1$ at \[
\Big(\frac{p}{1-q},0\Big) \qquad\text{and}\qquad \Big(\frac{p-q}{1-q},1\Big),\]
and these are the two endpoints of $\mathcal S$. Therefore
\[H_1(\Ber(p)\|\Ber(q))=\min\left\{
(1-q)\,h_{\mathsf{b}}\!\left(\frac{p}{1-q}\right),\,(1-q)\,h_{\mathsf{b}}\!\left(\frac{p-q}{1-q}\right)\right\}.\]
Using the symmetry $h_{\mathsf{b}}(u)=h_{\mathsf{b}}(1-u)$, the first term can be rewritten as
\[
(1-q)\,h_{\mathsf{b}}\!\left(\frac{1-p-q}{1-q}\right).
\]
Now
\[
0\le \frac{p-q}{1-q}\le \frac{1-p-q}{1-q}\le \frac12,
\]
because $q\le p\le \frac12$. Since $h_{\mathsf{b}}$ is increasing on $[0,\frac12]$, it follows that
\[
h_{\mathsf{b}}\!\left(\frac{p-q}{1-q}\right)
\le
h_{\mathsf{b}}\!\left(\frac{1-p-q}{1-q}\right)
=
h_{\mathsf{b}}\!\left(\frac{p}{1-q}\right).
\]
Hence
\[
H_1(\Ber(p)\|\Ber(q))
=
(1-q)\,h_{\mathsf{b}}\!\left(\frac{p-q}{1-q}\right)
\qquad\text{when } p>q.
\]

Combining the two cases, for all $0\le p,q\le \frac12$,
\[
H_1(\Ber(p)\|\Ber(q))
=
\begin{cases}
q\,h_{\mathsf{b}}\!\left(\dfrac{p}{q}\right), & p\le q,\\[1.2ex]
(1-q)\,h_{\mathsf{b}}\!\left(\dfrac{p-q}{1-q}\right), & p>q.
\end{cases}
\]
\end{proof}

\end{document}
